% ID: SEL-D0043-OB-E03
% ESTADO: cotejada. Segunda pasada visual y borrador compilado revisados.
% Fuente: Selectividad/sel_2014_matematicas_aplicadas/EXAMEN -ANDAL-2.pdf
% SHA-256: fc025a25faa338cbb9d8d288443f0322dfbcc0464f7474a16a84e8dfe47f81d1
% Página: 2; opción B; ejercicio original 3.

Se elige un número, al azar, entre el siguiente conjunto:
\[
\{225,201,162,210,180,172,156,193,218,167,176,222,215,120,190,171\}.
\]
\begin{enumerate}[label=\alph*)]
\item \textbf{(0.5 puntos)} Calcule la probabilidad de que el número elegido sea impar.
\item \textbf{(0.75 puntos)} Si el número elegido es múltiplo de 5, \textquestiondown{}cuál es la probabilidad de que sea mayor que 200?
\item \textbf{(0.75 puntos)} Determine si son independientes los sucesos $S$: ``el número elegido es mayor que 200'' y $T$: ``el número elegido es par''.
\item \textbf{(0.5 puntos)} Halle la probabilidad del suceso $S\cup T$.
\end{enumerate}
